\begin{slikaDesno}{fig/cap_mic.pdf}
\PID
На слици је приказан поједностављен модел капацитивног микрофона. Основни елемент функције микрофона је раван плочасти кондензатор, чија је једна електрода фиксирана за кућиште, док је друга причвршћена за еластичну мембрану која се деформише услед променљивог ваздушног притиска. Активна површина плоча износи $S = 4\unit{cm^2}$, док је равнотежно растојање између плоча (у одсуству деформације мембране) $d_0 = 100\unit{\upmu m}$.
Кондензатор је прикључен на идеалан извор сталне електромоторне силе $E = 42\unit{V}$ 
помоћу отпорника електричне отпорности $R = 1\unit{M\Omega}$. 
Микрофон је побуђен звучним таласом који изазива хармонијско кретање покретне електроде, тако да се растојање између плоча мења по закону $d(t) = d_0 + d_{\rm m}\sin(\omega t)$, где је амплитуда осцилација $d_{\rm m} = 10\unit{\upmu m}$, 
а кружна учестаност је 
$\upomega =10\unit{krad/s}$. Одредити амплитуду и почетну фазу струје $i$ успостављене у колу.
\end{slikaDesno}
\REZULTAT

Амплитуда тражене струје је $i_{\rm m} \approx 1,40\unit{\upmu A}$, а њена почетна фаза је 
$\uppsi
 \approx 109,5^\circ$.